\section*{LAMPIRAN}
\phantomsection
\addcontentsline{toc}{section}{LAMPIRAN}

\subsection*{Lampiran 1: Tautan Repositori dan Bukti Rilis}
\begin{itemize}
    \item \textbf{Repositori Aplikasi Android:} \url{https://github.com/coneckin/coneckin}
    \item \textbf{Landing Page Panduan:} \url{https://github.com/coneckin/coneckin.github.io}
    \item \textbf{Paket Rilis APK Terbaru:} \url{https://github.com/coneckin/coneckin/releases/latest/download/coneckin.apk}
\end{itemize}

\subsection*{Lampiran 2: Konfigurasi Pipeline CI/CD GitHub Actions}
Berikut adalah kutipan berkas konfigurasi alur kerja rilis otomatis (\texttt{.github/workflows/release.yml}) yang menandatangani biner secara otomatis dan mempublikasikan release APK/AAB:

\begin{verbatim}
name: Build and Release
on:
  push:
    branches: [main]
jobs:
  release:
    runs-on: ubuntu-latest
    steps:
      - uses: actions/checkout@v4
      - uses: actions/setup-java@v4
        with:
          distribution: temurin
          java-version: 17
      - uses: gradle/actions/setup-gradle@v4
      - name: Build APK dan AAB
        run: ./gradlew assembleRelease bundleRelease
      - name: Buat Release
        uses: softprops/action-gh-release@v2
        with:
          tag_name: v1.0.${{ github.run_number }}
          files: |
            coneckin.apk
            coneckin.aab
\end{verbatim}
